\documentclass[12pt]{article}

\usepackage[margin = .8in]{geometry}
\usepackage{amsmath}
\usepackage{graphicx}
\usepackage{multicol, enumerate, tabularx}

\usepackage[parfill]{parskip}

\usepackage{adjustbox, soul}

\usepackage{fancyhdr}
\pagestyle{fancy}

\lhead{\large\sf\textbf{MATH F113X: Finance}}
\chead{\large\sf\textbf{lecture notes}}
\rhead{\large\sf\textbf{Loans with Regular Payments}}

\renewcommand{\emph}[1]{\textsf{\textbf{#1}}}
\usepackage{tikz}
\usetikzlibrary{calc,trees,positioning,arrows,fit,shapes,through, backgrounds}
\usetikzlibrary{patterns}

\usetikzlibrary{decorations.markings}
\usetikzlibrary{arrows}

\usepackage{pgfplots}

\usepackage{longtable}
\usepackage{tabularx}

\newcommand{\ds}{\displaystyle}
\newcommand{\ans}[1][1in]{\rule{#1}{.5pt}}

\newcommand{\points}[1]{(#1 points.)}		% Trying to be lazy.

\usepackage{array}
\newcolumntype{L}[1]{>{\raggedright\let\newline\\\arraybackslash\hspace{0pt}}m{#1}}
\newcolumntype{C}[1]{>{\centering\let\newline\\\arraybackslash\hspace{0pt}}m{#1}}
\newcolumntype{R}[1]{>{\raggedleft\let\newline\\\arraybackslash\hspace{0pt}}m{#1}}
\newcommand{\red}[1]{\textcolor{red}{#1}}

\newcommand{\be}{\begin{enumerate}}
\newcommand{\ee}{\end{enumerate}}

\begin{document}
%\begin{center}
%{\large  Finance Section 4: Introduction to Loan Calculation}
%\end{center}
Goal: Think about loan calculations intuitively. What happens if you make regular payments?

\begin{enumerate}
\item  (Basic Example:) Suppose you take out a loan for \$10,000 with an annual interest rate of 5\% compounded \emph{annually}, and you will make payments annually, at the end of the year.
	\begin{enumerate}
	\item How much would you be charged in interest in year 1 if you didn't pay anything? What does that say about how much your monthly payment should be? (Round your answer to the nearest dollar.)
	\vfill
	\item Suppose you decided to make annual payments of \$600. How much do you owe at the end of one year? (One year of interest is added and one year of payments is subtracted, so $A = P + I - d$.)
	\vfill
	\item How much do you owe at the end of the second year? The third year?
	\vfill
	\item Use a \emph{spreadsheet} (fill down) to determine how long it will take to pay off the loan:
	
	{\tt
%	\hspace{-1cm}\begin{tabularx}{\textwidth}{|l||X|X|X|X|X|X|}
%	\hline
%	 & \verb'A' 		& \verb'B' 	& \verb'C' 		& \verb'D' & \verb'E' & \verb'F'\\ \hline \hline
%	1 & Interest Rate 	& Principal &  Year 		& Payment & Interest 		& Total \\ \hline
%	2 &	r 				&P			&t				& d 		&I				&A		\\ \hline		
%	2 & 0.05 			& 10000 	& 0 			& 0 		& 0				& =B2  \\ \hline
%	3 & 				& 			&\verb`=C2+1` 	& 600 		&\verb'= $A$2*$F$2' &  \verb'=E3+F2-D3' \\ \hline 
%	\end{tabularx}
	\hspace{-1cm}\begin{tabularx}{\textwidth}{|l||X|X|X|X|X|X|}
	\hline
	 & \verb'A' 		& \verb'B' 	& \verb'C' 		& \verb'D' & \verb'E' & \verb'F'\\ \hline \hline
	1 & Interest Rate 	& Principal &  payment 		& year & Interest 		& Total \\ \hline
	2 &	r 				&P			& d				& t 		&I				&A		\\ \hline		
	3 & 0.05 			& 10000 	& 600 			& 0 		& 0				& =B3  \\ \hline
	4 & 				& 			& 				& \verb`=C3+1` 		&\verb'= $A$3*F3' &  \verb'=E4+F3-$C$3' \\ \hline 
	\end{tabularx}


	}
	
	Number of years \ans\ Final payment \ans
	
	\item How much did you pay in total? How much was interest?
	\vfill
	\vfill
	\item What happens if you \emph{increase} your annual payments to \$800?
	\vfill
	\item What happens if you \emph{decrease} your annual payments to \$300?
	\vfill
	\end{enumerate}
\newpage
	
	\item \emph{Formulas:}
	
	\begin{tabular}{| l | l |}
\hline
$P$ = principal / starting amount & $r$ = annual interest rate (APR) \\ \hline
$I$ = interest & $n$ = number of compounding periods per year\\ \hline
$A$ = final amount & $t$ = length of the loan, in years\\ \hline
& $d$ = regular loan payment \\ \hline
\end{tabular}


\begin{center}
\begin{tabularx}{.8\linewidth}{|X | X|} 
\hline
Loan amount given payment & payment given loan amount\\ \hline
$\displaystyle P = \frac{d\left(1 - \left(1+\frac{r}{n}\right)^{(-nt)}\right)}{\left(\frac{r}{n}\right)}$ & $\displaystyle d = \frac{P\left(\frac{r}{n}\right)}{\left(1 - \left(1+\frac{r}{n}\right)^{(-nt)}\right)}$\\
\hline
\end{tabularx}
\end{center}


\item Suppose you take out a loan for \$10,000 with an annual interest rate of 5\% compounded \emph{monthly} and you are going to make payments monthly. If you want to pay off the loan in 3 years, what should your \emph{monthly payment} be? 

What are we looking for? \ans
\vfill


\item If you can afford a \$300 a month payment for 5 years, at an annual interest rate of 6.75\% compounded monthly, how much can you \emph{borrow}?   

What are we looking for? \ans

\vfill
	
	



\end{enumerate}
\end{document}
